\section{两个高斯分布的常用结论}

\subsection{独立高斯之和仍是高斯}
\label{app:gauss-sum}

\begin{proposition}
设 $\beps_1\sim\Ncal(\bzero,\sigma_1^2\bI)$、$\beps_2\sim\Ncal(\bzero,\sigma_2^2\bI)$ 相互独立，
$a,b\in\mathbb{R}$，则
\begin{equation}
  a\,\beps_1+b\,\beps_2\ \sim\ \Ncal\bigl(\bzero,\ (a^2\sigma_1^2+b^2\sigma_2^2)\bI\bigr).
  \label{eq:gauss-sum}
\end{equation}
\end{proposition}

\begin{pf}
设 $\bm{t}\in\mathbb{R}^d$，用特征函数。由于 $\beps_1,\beps_2$ 独立，
\begin{equation}
  \E\bigl[e^{\mathrm{i}\bm{t}^{\top}(a\beps_1+b\beps_2)}\bigr]
  =\E\bigl[e^{\mathrm{i}a\bm{t}^{\top}\beps_1}\bigr]\cdot\E\bigl[e^{\mathrm{i}b\bm{t}^{\top}\beps_2}\bigr]
  =e^{-\frac{1}{2}a^2\sigma_1^2\|\bm{t}\|^2}\cdot e^{-\frac{1}{2}b^2\sigma_2^2\|\bm{t}\|^2}
  =e^{-\frac{1}{2}(a^2\sigma_1^2+b^2\sigma_2^2)\|\bm{t}\|^2},
\end{equation}
这正是零均值高斯 $\Ncal(\bzero,(a^2\sigma_1^2+b^2\sigma_2^2)\bI)$ 的特征函数，
由特征函数与分布的一一对应即得结论。
\end{pf}

这个结论在正文中被用了两次：一次是在前向闭式解的归纳法里
（式~\eqref{eq:forward-induction}），把两步噪声合并成一步；
另一次是在推导 $\tilde{\bmu}_t$ 时把“关于 $\bx_{t-1}$ 的二次型”配方成高斯。

\subsection{两个高斯之间的 KL 散度}
\label{app:kl-gauss}

\begin{proposition}
设 $\bmu_1,\bmu_2\in\mathbb{R}^d$，$\bm{\Sigma}_1,\bm{\Sigma}_2$ 为 $d\times d$ 正定矩阵，则
\begin{equation}
  \begin{aligned}
  \Dkl\bigl(\Ncal(\bmu_1,\bm{\Sigma}_1)\,\big\|\,\Ncal(\bmu_2,\bm{\Sigma}_2)\bigr)
  =\frac{1}{2}\Bigl[&\operatorname{tr}\bigl(\bm{\Sigma}_2^{-1}\bm{\Sigma}_1\bigr)
   +(\bmu_2-\bmu_1)^{\top}\bm{\Sigma}_2^{-1}(\bmu_2-\bmu_1) \\
   &-d+\ln\frac{\det\bm{\Sigma}_2}{\det\bm{\Sigma}_1}\Bigr].
  \end{aligned}
  \label{eq:kl-gauss}
\end{equation}
特别地，当两者协方差相同（$\bm{\Sigma}_1=\bm{\Sigma}_2=\sigma^2\bI$）时，
只有均值不同，公式退化成
\begin{equation}
  \Dkl\bigl(\Ncal(\bmu_1,\sigma^2\bI)\,\big\|\,\Ncal(\bmu_2,\sigma^2\bI)\bigr)
  =\frac{1}{2\sigma^2}\bigl\|\bmu_1-\bmu_2\bigr\|^2 .
  \label{eq:kl-gauss-equal-var}
\end{equation}
\end{proposition}

式~\eqref{eq:kl-gauss-equal-var} 就是正文式~\eqref{eq:Lt-kl} 的来源：
它把“两个分布有多远”这件事，换成了“两个均值差多少”这件更简单的事，
而后者正好可以用平方误差来优化。

\subsection{ELBO 三项分解的完整展开}
\label{app:elbo}

从式~\eqref{eq:ddpm-elbo} 出发，把两个链式乘积代入：
\begin{equation}
  \begin{aligned}
    \mathcal{L}(\bx_0)
    &=\E_q\Bigl[\log p(\bx_T)+\sum_{t=1}^{T}\log p_\theta(\bx_{t-1}\mid\bx_t)
      -\sum_{t=1}^{T}\log q(\bx_t\mid\bx_{t-1})\Bigr] \\
    &=\E_q\Bigl[\log p(\bx_T)+\log p_\theta(\bx_0\mid\bx_1) \\
    &\qquad\qquad+\sum_{t=2}^{T}\log p_\theta(\bx_{t-1}\mid\bx_t)
      -\sum_{t=2}^{T}\log q(\bx_t\mid\bx_{t-1})-\log q(\bx_1\mid\bx_0)\Bigr].
  \end{aligned}
  \label{eq:app-elbo-1}
\end{equation}
对 $t\ge2$ 的每一对相邻项用贝叶斯公式~\eqref{eq:bayes-trick}：
\begin{equation}
  \begin{aligned}
    \sum_{t=2}^{T}\log q(\bx_t\mid\bx_{t-1})
    &=\sum_{t=2}^{T}\log q(\bx_{t-1}\mid\bx_t,\bx_0) \\
    &\quad+\underbrace{\sum_{t=2}^{T}\bigl[\log q(\bx_t\mid\bx_0)-\log q(\bx_{t-1}\mid\bx_0)\bigr]}_{\text{相消}},
  \end{aligned}
  \label{eq:app-elbo-2}
\end{equation}
后一个求和是相邻两项的差，只剩首尾：$\log q(\bx_T\mid\bx_0)-\log q(\bx_1\mid\bx_0)$。
代回式~\eqref{eq:app-elbo-1}，其中 $-\log q(\bx_1\mid\bx_0)$ 恰好被抵消，于是
\begin{equation}
  \mathcal{L}(\bx_0)
  =\E_q\Bigl[\log p_\theta(\bx_0\mid\bx_1)
   -\sum_{t=2}^{T}\log\frac{q(\bx_{t-1}\mid\bx_t,\bx_0)}{p_\theta(\bx_{t-1}\mid\bx_t)}
   -\log\frac{q(\bx_T\mid\bx_0)}{p(\bx_T)}\Bigr].
  \label{eq:app-elbo-3}
\end{equation}
现在逐项看：对固定的 $t$，第二项里除 $\bx_t$ 之外的变量都能积掉，期望退化为
$\E_{q(\bx_t\mid\bx_0)}\bigl[\Dkl(q(\bx_{t-1}\mid\bx_t,\bx_0)\|p_\theta(\bx_{t-1}\mid\bx_t))\bigr]$；
第三项就是 $\Dkl(q(\bx_T\mid\bx_0)\|p(\bx_T))$；第一项是 $-\log p_\theta(\bx_0\mid\bx_1)$ 的期望。于是
\begin{equation}
  \begin{aligned}
    \mathcal{L}(\bx_0)
    &=-\underbrace{\Dkl\bigl(q(\bx_T\mid\bx_0)\|p(\bx_T)\bigr)}_{L_T} \\
    &\quad-\sum_{t=2}^{T}\underbrace{\E_{q(\bx_t\mid\bx_0)}
      \bigl[\Dkl\bigl(q(\bx_{t-1}\mid\bx_t,\bx_0)\|p_\theta(\bx_{t-1}\mid\bx_t)\bigr)\bigr]}_{L_{t-1}} \\
    &\quad-\underbrace{\E\bigl[-\log p_\theta(\bx_0\mid\bx_1)\bigr]}_{L_0},
  \end{aligned}
  \label{eq:app-elbo-4}
\end{equation}
即 $\mathcal{L}(\bx_0)=-(L_T+\sum_{t=2}^{T}L_{t-1}+L_0)$，与正文式~\eqref{eq:elbo-three} 一致。
